\documentclass[11pt]{article}
\usepackage{enumitem}
\usepackage{geometry}
\usepackage{xcolor}
\usepackage[normalem]{ulem}
\geometry{margin=1in}

% --- 2026 revision markup -------------------------------------------
% \was{...}  prior wording, kept inline in gray beneath the revision
% \newprof   flags a proficiency added in this revision
% \cut{...}  a proficiency removed in this revision, kept for the record
\newcommand{\was}[1]{\newline{\small\textcolor{gray}{\textit{was:} #1}}}
\newcommand{\newprof}{\ {\small\textcolor{gray}{(new)}}}
\newcommand{\cut}[1]{\item[\textcolor{gray}{\small$\times$}]%
  \textcolor{gray}{\sout{#1}\ \textit{(removed)}}}
% \priorunits{...}  prior unit-level structure, kept for the record
\newcommand{\priorunits}[1]{\par\medskip
  {\small\textcolor{gray}{\textit{Prior unit structure:} #1}}}
% \revnote{...}  editorial note about this revision
\newcommand{\revnote}[1]{\par\medskip
  {\small\textcolor{gray}{\textit{Note:} #1}}}
% \prof{CODE}{text}  coded proficiency (matches the canonical calculus document)
\newcommand{\prof}[2]{\item[\textbf{#1}]#2}
% \unithead{title}{days}  unit banner inside a coded proficiency list
\newcommand{\unithead}[2]{\par\medskip\noindent
  \textbf{#1}\hfill{\small\textit{#2}}\par\smallskip}

\title{NWA Mathematics Proficiencies}
\date{\today}

\begin{document}
\maketitle

\section{Pre Algebra 1 (6th Grade)}

\begin{enumerate}
\item Understand and apply basic geometric concepts
\item Calculate perimeter and area of basic shapes
\item Understand and calculate the volume of a prism\newprof
\item Understand and work with ratios
\item Compare and order fractions
\item Add and subtract fractions with like denominators
\item Understand decimal place value
\item Add and subtract decimals
\item Multiply decimals
\item Divide decimals
\item Convert between fractions and decimals
\item Understand and use percents
\item Convert between fractions decimals and percents
\item Identify and extend patterns
\item Write and evaluate algebraic expressions
\item Understand the concept of variables
\item Use the order of operations
\item Understand and use the distributive property
\item Combine like terms
\item Understand coordinate plane basics
\item Plot points on a coordinate plane
\item Understand and identify types of angles
\item Measure angles with a protractor
\item Understand properties of triangles
\item Solve one-step equations
\item Understand and use inequalities
\item Collect and organize data
\item Create and interpret bar graphs
\item Create and interpret line graphs
\item Find mean median and mode
\item Understand basic probability concepts
\item Calculate simple probabilities
\item Use tree diagrams for counting
\item Understand prime and composite numbers
\end{enumerate}

\section{Pre Algebra 2 (7th Grade)}

\begin{enumerate}
\item Add and subtract fractions with unlike denominators
\item Understand and work with proportional relationships
\item Solve proportion problems
\item Calculate area and perimeter of complex shapes
\item Understand and use scale drawings
\item Work with similar figures
\item Understand and use negative numbers
\item Add and subtract integers
\item Find greatest common factor and least common multiple
\item Simplify fractions
\item Multiply and divide fractions
\item Work with mixed numbers
\item Understand and use scientific notation
\item Multiply and divide in scientific notation
\item Understand and work with rational numbers
\item Compare and order rational numbers
\item Understand absolute value
\item Graph inequalities on a number line
\item Write and solve inequalities
\item Understand and use rates and unit rates
\item Solve multi-step equations
\item Solve multi-step inequalities
\item Factor linear expressions
\item Expand linear expressions
\item Determine experimental probability
\item Determine theoretical probability
\item Use probability to make predictions
\end{enumerate}

\section{Pre Algebra 3 (8th Grade)}

\begin{enumerate}
\item Understand and perform operations with rational numbers
\item Simplify expressions with rational numbers
\item Evaluate expressions with rational numbers
\item Distinguish between and compare the value of rational and irrational
      numbers\newprof
\item Understand and use properties of exponents
\item Apply laws of exponents
\item Work with powers of 10 and scientific notation
\item Perform operations in scientific notation
\item Apply the Pythagorean theorem
\item Understand and perform geometric transformations
\item Graph transformations on the coordinate plane
\item Understand similarity and congruence
\item Understand and apply properties of similarity to find unknown angles and side
      lengths\newprof
\item Identify and use angle relationships
\item Understand the concept of a function
\item Determine if a relation is a function
\cut{Use function notation}
\item Evaluate functions
\item Graph linear functions
\item Find slope from a graph or two points
\item Write equations of lines
\item Use linear equations to model real-world scenarios\newprof
\item Identify parallel and perpendicular lines
\item Solve systems of linear equations graphically
\item Solve systems of linear equations algebraically
\item Create and interpret scatter plots
\item Identify associations in categorical data\newprof
\item Understand correlation and causation
\end{enumerate}
\revnote{Reviewed by Dan, August 2026. Items 21--28 of the prior list were an exact
  duplicate of items 13--20 (concept of a function, determine if a relation is a
  function, use function notation, evaluate functions, graph linear functions, find
  slope, write equations of lines, identify parallel and perpendicular lines) and have
  been deleted as a copy-paste error rather than a curricular change.}

\section{Integrated Math 1 (9th Grade)}

\begin{enumerate}
\item Identify domain and range\newprof
\item Use interval notation\newprof
\item Use function notation
      \was{Understand and use function notation}
\item Identify continuous vs.\ discrete data\newprof
\item Write explicit equations\newprof
\item Write recursive equations\newprof
\item Graph a function from a table or equation
      \was{Graph and analyze linear functions}
\item Write exponential equations\newprof
\item Write linear equations
      \was{Write equations of lines in various forms}
\item Understand the difference between linear and exponential functions\newprof
\item Understand the relationship between fractional exponents and radicals\newprof
\item Convert negative exponential expressions into positive exponential
      expressions\newprof
\item Solve a multi-step equation\newprof
\item Solve a multi-step inequality\newprof
\item Rearrange equations to isolate a variable\newprof
\item Graph inequalities on the coordinate plane\newprof
\item Graph inequalities on a number line\newprof
\item Do transformations using construction\newprof
\item Do transformations on a coordinate plane\newprof
\item Graph parallel and perpendicular lines
      \was{Identify parallel and perpendicular lines}
\item Construct parallel and perpendicular lines\newprof
\item Understand rotational and reflectional properties of quadrilaterals\newprof
\item Solve a system of linear equations using substitution
      \was{Solve systems of linear equations}
\item Solve a system of linear equations using elimination\newprof
\item Solve a system of linear equations by graphing\newprof
\item Interpret box and whisker plots\newprof
\item Create a histogram\newprof
\item Interpret a histogram\newprof
\item Understand standard deviation as a measure of spread\newprof
\cut{Analyze bivariate data and correlation}
\end{enumerate}
\priorunits{Unit 1: Sequences \textbullet\ Unit 2: Linear and Exponential
  Functions \textbullet\ Unit 3: Features of Functions \textbullet\ Unit 4:
  Equations and Inequalities \textbullet\ Unit 5: Systems of Equations and
  Inequalities \textbullet\ Unit 6: Two-way Tables \textbullet\ Unit 7: Functions
  and Their Graphs \textbullet\ Unit 8: Associations in Data \textbullet\ Unit 9:
  Congruence and Rigid Transformations}

\section{Integrated Math 2 (10th Grade)}

\begin{enumerate}
\item Unit 1: Quadratic Functions and Their Model Patterns and Real-World Contexts
      \was{Unit 1: Functions and Their Inverses}
\item Unit 2: Structures of Quadratic Functions (forms and graph features)
      \was{Unit 2: Logarithmic Functions}
\item Unit 3: Solving Quadratic Functions (using graphs and completing the square,
      the quadratic formula, and examining strategies)
      \was{Unit 3: Polynomial Functions}
\item Unit 4: Piecewise, Absolute Value, and Inverse Equations
      \was{Unit 4: Polynomial Equations}
\item Unit 5: Rational Exponents and Radicals
\item Unit 6: Similarity and Right Triangle Trigonometry
\item Unit 7: Circles
\item Unit 8: Conditional Probability
\item Unit 9: Modeling Data
\item Unit 10: Constructions and Proof
\end{enumerate}

\section{Integrated Math 3 (11th Grade)}

\begin{enumerate}[itemsep=6pt]

\item \textbf{Unit 1: Functions and Their Inverses}
      \was{Unit 1: Functions and Polynomials}
  \begin{itemize}[nosep,leftmargin=1.5em]
  \item Create an inverse function from a story context by changing the input and
        output variables
  \item Recognize inverse functions in tables, graphs, and equations
  \item Determine if a function is invertible
  \item Understand that inverse functions ``undo'' each other
  \end{itemize}

\item \textbf{Unit 2: Logarithms}
      \was{Unit 2: Exponential and Logarithmic Functions}
  \begin{itemize}[nosep,leftmargin=1.5em]
  \item Find the exact value of a logarithmic expression when the argument of the
        logarithm is a power of the base
  \item Estimate the value of a logarithmic expression when an exact value cannot be
        found
  \item Recognize and use the properties of logarithms
  \item Graph logarithmic functions
  \item Use the properties of logarithms to find equivalent expressions
  \item Use base 10 (common) logarithms to solve exponential equations
  \item Solve logarithmic equations in base 10
  \item Understand compound interest, including the continuous compound interest
        formula
  \item Understand how different factors (interest rate, compounding periods, number
        of years invested) affect the amount earned from compound interest
  \item Identify features of a base $e$ exponential function
  \item Model continuous growth and decay using base $e$ exponential functions
  \item Solve exponential equations using common logarithms (base 10) and natural
        logarithms (base $e$)
  \end{itemize}

\item \textbf{Unit 3: Polynomials}
      \was{Unit 3: Rational Functions}
  \begin{itemize}[nosep,leftmargin=1.5em]
  \item Identify a cubic function from a rate of change
  \item Describe the features of cubic functions
  \item Graph cubic functions
  \item Describe similarities and differences between cubic functions and quadratic
        functions
  \item Add and subtract polynomials both algebraically and geometrically
  \item Multiply polynomials using the distributive property
  \item Use Pascal's Triangle to raise a binomial to a power
  \item Divide polynomials using long division
  \item Write two equivalent polynomial multiplication statements for each division
        statement
  \item Understand the Polynomial Remainder Theorem
  \item Use the Fundamental Theorem of Algebra to determine the number of roots a
        polynomial has
  \item Explain why polynomials are closed under addition, subtraction, and
        multiplication, but not closed under division
  \item Write a polynomial in its factored form, given the roots of the polynomial
  \item Find the other roots of a polynomial, given a factor or not
  \item Describe pairs of irrational or imaginary roots of polynomials
  \item Determine the end behavior of a polynomial of a given degree
  \item Use proper notation to write the end behavior of a given polynomial
  \item Graph a polynomial given information about its roots and degree
  \item Determine if a function is even or odd
  \item Write the equation of a polynomial in standard and factored form
  \end{itemize}

\item \textbf{Unit 4: Polynomial Identities and Proofs}
      \revnote{No revised proficiencies supplied for this unit; the prior unit title is
        retained pending department input.}

\item \textbf{Unit 5: Modeling with Geometry}
      \was{Unit 5: Modeling with Functions}
  \begin{itemize}[nosep,leftmargin=1.5em]
  \item Visualize, describe, or draw the cross-sections of a 3D solid intersected by
        a plane
  \item Visualize, describe, or draw solids of revolution formed when a 2D figure is
        rotated about an axis
  \item Approximate the volume of solids of revolution using formulas for area and
        volume by decomposing the solid into smaller pieces
  \item Explain the concept of density
  \item Find the exact measures of missing sides and angles in special right triangles
  \item Use the Law of Cosines
  \item Use the Law of Sines
  \item Identify when to use the Law of Sines or the Law of Cosines
  \end{itemize}

\item \textbf{Unit 6: Modeling Periodic Behavior}
      \was{Unit 6: Trigonometric Functions}
  \begin{itemize}[nosep,leftmargin=1.5em]
  \item Apply right triangle trigonometry to calculate how far points on the circle
        are above or below the center of the circle, or how far to the right or left
        of the center of the circle
  \item Calculate the angular speed given the period of rotation of a Ferris wheel,
        and calculate the period of rotation of a Ferris wheel given its angular speed
  \item Identify the amplitude, period, and location of the midline from a
        trigonometric function
  \item Sketch the graph of a function given the amplitude, period, and midline
  \item Calculate arc length for angles of rotation measured in degrees or radians,
        and find the angle of rotation in radians or degrees given the arc length
  \item Calculate the sine and cosine of an angle of rotation on a circle with a given
        radius of $r$
  \item Calculate the sine and cosine of an angle of rotation on a unit circle
  \item Label a point on the circle with rectangular coordinates or polar coordinates
  \item Write equations to model the vertical or horizontal positions of a point
        rotating around a circle using either degrees or radians
  \end{itemize}

\item \textbf{Unit 7: Trigonometric Functions, Equations, and Identities}
      \was{Unit 7: Trigonometric Equations and Identities}
  \begin{itemize}[nosep,leftmargin=1.5em]
  \item Write trigonometric functions that represent horizontal transformations of the
        graph or different initial starting positions of the circular motion
  \item Represent the same trigonometric graph using either sine or cosine functions
        and using either degrees or radians to measure angles of rotation
  \item Write trigonometric functions to model periodic behavior that is not about
        circular motion
  \item Use the inverse of trigonometric functions represented with tables, graphs, or
        equations as needed
  \item Define the tangent function in terms of angles of rotation and sketch its graph
  \item Find tangent values on the unit circle
  \item Understand and explain the fundamental trigonometric identities
  \item Use trigonometric identities and inverse trigonometric functions to solve
        trigonometric equations
  \item Understand and explain the derivation of the sum and difference identities and
        the double angle identities for sine, cosine, and tangent
  \item Understand and use the restrictions on the definitions of the inverse
        trigonometric functions
  \item Use algebra, trigonometric identities, and the inverse trigonometric functions
        to find all solutions of trigonometric functions
  \end{itemize}

\item \textbf{Unit 8: Modeling with Functions}
      \was{Unit 8: Statistics}
  \begin{itemize}[nosep,leftmargin=1.5em]
  \item Create a table or a graph for a transformation of a parent function whether
        the transformation is described in geometric notation
  \item Describe the effects of a transformation on a graph
  \item Determine the effects of combining functions from two different function
        families using addition or multiplication
  \item Find the composition of two functions defined by tables
  \item Find the function rule for the composition of two or more functions defined
        with algebraic rules
  \item Decompose a function into its component functions
  \end{itemize}

\item \textbf{Unit 9: Statistics}
      \was{Unit 9: Optimization}
  \begin{itemize}[nosep,leftmargin=1.5em]
  \item Describe the features of a normal distribution
  \item Know and use the 68--95--99.7 rule to interpret a point in a distribution
  \item Calculate and interpret a z-score
  \item Estimate population percentages using a normal distribution as a model and
        using technology or a table to help estimate
  \item Determine whether or not a sample is random
  \item Identify different sample types and describe their purposes
  \item Determine the population, sample, and parameter of interest in a study
  \item Identify the differences between an experimental study, an observational
        study, and a survey and describe the strengths and weaknesses of each type
  \item Use a simulation to predict the likelihood of a given event
  \item Develop a plan to study a parameter of interest using an appropriate sampling
        strategy and study design
  \end{itemize}

\end{enumerate}
\revnote{Unit-level proficiencies supplied by the IM3 instructor, August 2026,
  converted from ``I can'' statements to the imperative voice used throughout this
  document. Wording is otherwise unchanged.}

\section{Pre-Calculus (12th Grade)}

\begin{enumerate}
\item Graph exponential and logarithmic functions
\item Solve exponential and logarithmic equations
\item Use properties of logarithms
\item Model with exponential and logarithmic functions
\item Understand and work with complex numbers
\item Perform operations with complex numbers
\item Find complex zeros of polynomials
\item Use the Fundamental Theorem of Algebra
\item Graph polynomial functions
\item Analyze rational functions and asymptotes
\item Perform polynomial and rational inequalities
\item Graph and analyze conic sections
\item Define and evaluate trigonometric functions
\item Graph trigonometric functions
\item Transform trigonometric functions
\item Evaluate inverse trigonometric functions
\item Verify trigonometric identities
\item Solve trigonometric equations
\item Apply sum and difference formulas
\item Use double and half-angle formulas
\item Understand and apply Laws of Sines and Cosines
\item Work with vectors in the plane
\item Perform vector operations
\item Apply vectors to solve problems
\item Work with parametric equations
\item Graph parametric equations
\item Convert between parametric and rectangular forms
\item Understand and work with sequences
\item Understand and work with series
\end{enumerate}

\section{Statistics 1 (12th Grade)}

\begin{enumerate}
\item Conceptualize and operationalize variables\newprof
\item Understand the iterative process of forming research questions and collecting
      data\newprof
\item Distinguish among different types of data\newprof
\item Understand the issues involved in cleaning up data\newprof
\item Describe and compare distributions using stemplots, histograms, box plots,
      and ntigrams
      \was{Describe and compare distributions}
\item Calculate measures of center and spread: mean, median, midrange, and mode;
      MAD, standard deviation, five-number summaries, and outlier tests
      \was{Calculate measures of center and spread}
\item Calculate and interpret z-scores\newprof
\item Use normal distributions, including percentiles
      \was{Use normal distributions}
\item Distinguish among percents, percent change, percentage point change, and
      percentiles\newprof
\item Calculate and interpret index scores\newprof
\item Understand sampling distributions
\item Understand sampling methods and their attendant concerns\newprof
\item Design surveys, experiments, and observational studies
      \was{Design surveys and experiments}
\item Understand bias and variability
\item Calculate and interpret correlation, including its non-transitivity and the
      importance of spurious associations and lurking variables
      \was{Calculate and interpret correlation}
\item Use scatter plots and the quadrant count ratio to describe association\newprof
\item Perform linear regression
\item Assess regression models
\item Transform data for linearity
\item Interpret models with more than one explanatory variable (multivariable
      regression)\newprof
\item Understand probability fundamentals, including the basic probability rules
      \was{Understand probability fundamentals}
\item Calculate conditional probability and apply Bayes' theorem
      \was{Calculate conditional probability}
\item Use probability distributions
\item Work with binomial distributions
\item Work with geometric distributions
\item Understand the Central Limit Theorem
\item Calculate confidence intervals for proportions
\item Calculate confidence intervals for means
\cut{Perform hypothesis tests for proportions}
\cut{Perform hypothesis tests for means}
\item Understand Type I and Type II errors, and interpret the confidence and power
      of an experiment
      \was{Understand Type I and Type II errors}
\item Calculate and interpret p-values, understand their meaning and limitations,
      and use confidence intervals to report results and quantify effect sizes
      \was{Calculate and interpret p-values}
\item Understand the false discovery rate and the false skepticism rate\newprof
\item Compare two proportions
\item Compare two means
\cut{Perform chi-square tests}
\cut{Test for independence}
\cut{Test for goodness of fit}
\item Analyze categorical data, including calculating and interpreting odds ratios
      \was{Analyze categorical data}
\item Create and interpret statistical reports
\end{enumerate}

\section{Statistics 2 (13th Grade)}

\begin{enumerate}
\item Distinguish between samples and populations
\item Classify different types of data
\item Apply appropriate sampling techniques
\item Analyze ethical considerations in statistical studies
\item Construct and interpret frequency tables
\item Create and analyze stem-and-leaf plots and histograms
\item Construct and interpret ogives pie charts bar charts and Pareto charts
\item Create stacked column charts and time series plots
\item Construct and analyze scatter plots for bivariate data
\item Calculate and interpret measures of center
\item Calculate and interpret measures of spread
\item Determine and interpret measures of position
\item Calculate and interpret correlation coefficients
\item Perform linear regression analysis and interpret results
\item Apply basic probability rules including complement rule
\item Calculate probabilities using union and intersection of events
\item Determine independence of events and apply conditional probability
\item Apply counting rules including permutations and combinations
\item Identify and work with probability distributions
\item Apply the binomial distribution to solve probability problems
\item Use the normal distribution to calculate probabilities and z-scores
\item Apply the central limit theorem to sampling distributions
\item Use the Student's t-distribution for appropriate statistical inference
\item Construct and interpret confidence intervals for a population proportion
\item Construct and interpret confidence intervals for a population mean
\item Conduct hypothesis tests for a single population proportion
\item Conduct hypothesis tests for a single population mean
\item Identify and explain Type I and Type II errors in hypothesis testing
\item Construct confidence intervals for comparing two population proportions
\item Conduct hypothesis tests comparing two population proportions
\item Construct confidence intervals for comparing two population means
\item Conduct hypothesis tests comparing two population means
\item Perform chi-squared goodness of fit tests
\item Conduct chi-squared tests for independence
\item Interpret results of chi-squared analyses
\item Understand and apply the F distribution
\item Conduct one-way ANOVA to compare multiple means
\item Perform pairwise comparison of means using appropriate methods
\item Perform comprehensive linear regression analysis
\item Conduct hypothesis tests for correlation coefficients
\item Conduct hypothesis tests for linear regression parameters
\item Calculate and interpret prediction intervals
\item Apply statistical methods to analyze real-world data sets
\item Present statistical findings clearly and accurately
\item Evaluate the validity and limitations of statistical conclusions
\end{enumerate}

\section{Advanced Pre-Calculus (11th-12th)}

\begin{enumerate}
\item Graph functions and their transformations
\item Solve equations graphically and algebraically
\item Graph exponential and logarithmic functions
\item Solve problems involving exponents and logarithms
\item Graph quadratic and higher degree polynomial functions
\item Find zeros of polynomial functions
\item Use long division and synthetic division of polynomials
\item Perform operations with complex numbers
\item Determine complex zeros using the Fundamental Theorem of Algebra
\item Graph rational functions
\item Find asymptotes of rational functions
\item Understand and correctly use both radian and degree measure
\item Define and evaluate trigonometric functions using the unit circle
\item Use reference angles to evaluate trigonometric functions
\item Graph trigonometric functions and their transformations
\item Evaluate inverse trigonometric functions
\item Define and evaluate trigonometric functions using right triangles
\item Use trigonometric formulas and laws of sines and cosines
\item Verify trigonometric identities using fundamental identities
\item Apply addition and subtraction formulas for trigonometric functions
\item Use double-angle and half-angle formulas
\item Solve trigonometric equations
\item Use limit laws to calculate limits as variable approaches constant
\item Demonstrate how to find the slope of a tangent line to a curve
\item Determine continuity of a function on a given domain
\item Calculate a derivative using the definition
\item Demonstrate understanding of the derivative as a function
\item Recall basic algebraic and trigonometric derivatives from memory
\item Calculate derivatives using product and quotient rules
\item Calculate derivatives using the chain rule
\item Calculate higher derivatives
\item Use implicit differentiation to find a derivative
\item Use derivatives to describe rates of change
\item Solve related rate problems
\item Use linear approximation to solve problems
\end{enumerate}

\section{Combined Calculus 1+2 (12th-13th)}

\noindent\textit{Dual credit: Portland State University MTH 251 (Calculus I) and
MTH 252 (Calculus II). Forty-five proficiencies across ninety-three class days.
Codes are stable identifiers, not positions --- the list runs in teaching order, so
the numbers do not.}
\revnote{Reconciled August 2026 against the canonical document
  \texttt{calculus\_proficiencies26-27.pdf}, in the \texttt{calc\_combo} admin
  proficiencies folder. That document governs; this section mirrors it.}

\unithead{Continuous Skills}{all year}
\begin{itemize}[leftmargin=3.4em,labelsep=0.5em,labelwidth=2.6em,align=left,itemsep=2pt]
\item[] Simplify expressions using algebra skills\newprof
\item[] Demonstrate an understanding of mathematical concepts and reasoning by
        appropriately using definitions, arguments, and proofs\newprof
\item[] Apply and adapt various strategies for problem solving at the calculus level,
        showing perseverance and creativity\newprof
\item[] Communicate mathematical thinking clearly to others, including written work
        that is thorough, organized, and carefully labeled, and oral explanations that
        are fluent\newprof
\item[] Complete and submit assignments on time\newprof
\item[] Use graphing software to investigate and verify functions and their graphs
        (a graphing calculator, GeoGebra or Desmos --- any is fine)
        \was{Use a graphing calculator}
\end{itemize}

\unithead{Unit 1 --- Review of Limits and Derivatives}{Days 1--15, Chapters 1--2}
\begin{itemize}[leftmargin=3.4em,labelsep=0.5em,labelwidth=2.6em,align=left,itemsep=2pt]
\prof{P02}{Describe basic classes of functions graphically, algebraically and verbally}
\prof{P03}{Transform functions using translation, scaling and reflection}
\prof{P04}{Use limit laws to calculate limits as the variable approaches a constant}
\prof{P06}{Determine the continuity of a function on a given domain}
\prof{P05}{Demonstrate how to find the slope of a tangent line to a curve}
\prof{P07}{Calculate a derivative using the definition}
\prof{P08}{Demonstrate understanding of the derivative as a function}
\prof{P09}{Recall the basic algebraic and trigonometric derivatives from memory}
\prof{P10}{Calculate appropriate derivatives using the product and quotient rules}
\prof{P12}{Calculate higher derivatives}
\prof{P11}{Calculate derivatives using the chain rule}
\prof{P13}{Use implicit differentiation to find a derivative}
\prof{P14}{Use derivatives to describe rates of change}
\prof{P15}{Solve related rate problems}
\prof{P16}{Use linear approximation to solve problems}
\prof{P43}{Evaluate Taylor polynomials}
\prof{P25}{Use L'H\^{o}pital's rule to determine limits of indeterminate forms
      \was{Use L'H\^{o}pital's rule to determine limits of quotients}}

\end{itemize}
\unithead{Unit 2 --- Applications of Differentiation}{Days 16--27, Chapter 3}
\begin{itemize}[leftmargin=3.4em,labelsep=0.5em,labelwidth=2.6em,align=left,itemsep=2pt]
\prof{P23}{Calculate critical points to determine the extreme values of a function}
\prof{P21}{Understand the Mean Value Theorem}
\prof{P24}{Sketch the graph of a function using first and second derivatives}
\prof{P22}{Use limit laws to calculate limits as the variable approaches infinity}
\prof{P26}{Solve optimization problems}
\prof{P44}{Use Newton's method to approximate the roots of an equation\newprof}
\prof{P27}{Find the antiderivative of a basic function}

\end{itemize}
\unithead{Unit 3 --- The Integral}{Days 28--39, Chapter 4}
\begin{itemize}[leftmargin=3.4em,labelsep=0.5em,labelwidth=2.6em,align=left,itemsep=2pt]
\prof{P28}{Use rectangular approximations and integrals to find area under curves}
\prof{P29}{Evaluate definite and indefinite integrals}
\prof{P32}{Apply the Fundamental Theorem of Calculus}
\prof{P30}{Solve problems involving net change}
\prof{P33}{Use integration methods: substitution \textit{(introduced here; assessed
      again in full at Unit 6)}}

\end{itemize}
\unithead{Unit 4 --- Applications of Integration}{Days 40--55, Chapter 5}
\begin{itemize}[leftmargin=3.4em,labelsep=0.5em,labelwidth=2.6em,align=left,itemsep=2pt]
\prof{P36}{Use integrals to determine the area between two curves}
\prof{P38}{Calculate the volume of a solid with given base and cross section}
\prof{P37}{Calculate volumes using disk, annulus and shell methods}
\prof{P31}{Calculate volumes of revolution using discs or cylindrical shells}
\prof{P45}{Calculate the average value of a function, and solve problems involving
      density\newprof}
\prof{P40}{Solve problems involving applications to physics and engineering}

\end{itemize}
\unithead{Unit 5 --- Inverse Functions}{Days 56--65, Chapter 6}
\begin{itemize}[leftmargin=3.4em,labelsep=0.5em,labelwidth=2.6em,align=left,itemsep=2pt]
\prof{P17}{Determine the inverse of a function when appropriate}
\prof{P18}{Use exponential and logarithmic laws to simplify expressions}
\prof{P19}{Calculate derivatives involving exponential and logarithmic functions}
\prof{P39}{Solve problems involving exponential growth and decay}
\prof{P20}{Calculate derivatives involving inverse trigonometric functions}

\end{itemize}
\unithead{Integration Toolkit}{Days 66--68, prerequisite repair for Chapter 7}
\begin{itemize}[leftmargin=3.4em,labelsep=0.5em,labelwidth=2.6em,align=left,itemsep=2pt]
\item[] No new proficiencies. Re-assesses P33 (substitution), polynomial division and
        trigonometric identities.

\end{itemize}
\unithead{Unit 6 --- Techniques of Integration}{Days 69--83, Chapter 7}
\begin{itemize}[leftmargin=3.4em,labelsep=0.5em,labelwidth=2.6em,align=left,itemsep=2pt]
\prof{P33}{Use integration methods: substitution, parts, trigonometric integrals,
      trigonometric substitution, and partial fractions
      \was{Use integration methods: substitution parts partial fractions}}
\prof{P34}{Use Simpson, trapezoidal and midpoint approximations}
\prof{P35}{Identify improper integrals and evaluate them using limits and comparisons
      \was{Identify improper integrals and know how to evaluate them}}

\end{itemize}
\unithead{Unit 7 --- Further Applications of Integration}{Days 84--93, Chapter 8}
\begin{itemize}[leftmargin=3.4em,labelsep=0.5em,labelwidth=2.6em,align=left,itemsep=2pt]
\prof{P42}{Calculate arc length and surface area}
\prof{P41}{Use integrals to determine probabilities}
\end{itemize}

\end{document}